\documentclass[11pt]{article}

\usepackage[margin=1in]{geometry}
\usepackage{amsmath, amssymb}
\usepackage{booktabs}
\usepackage{array}
\usepackage{xcolor}
\usepackage{listings}
\usepackage{enumitem}

\lstset{
  basicstyle=\ttfamily\small,
  breaklines=true,
  frame=single,
  columns=fullflexible
}

\title{Deriving an Efficient Distributed Computation of $$GG^\top$$ from Column-Sharded Gradients}
\author{}
\date{}

\begin{document}
\maketitle

\section{The problem we want to solve}

Suppose we are maintaining the left preconditioner in SOAP/Shampoo. At each update, we need the exponential moving average

$$
L \leftarrow \beta L + (1-\beta)GG^\top,
$$

where

$$
G \in \mathbb{R}^{m \times n}.
$$

The hard part is that, under tensor parallelism, $$G$$ is not stored on one GPU. Instead, it is column-sharded across $$p$$ ranks:

$$
G =
\begin{bmatrix}
G_0 & G_1 & \cdots & G_{p-1}
\end{bmatrix},
$$

where each rank $$r$$ owns

$$
G_r \in \mathbb{R}^{m \times n_{\text{local}}}.
$$

So the question is:

\begin{quote}
How do we compute and store $$GG^\top$$ efficiently when every rank only owns a column shard of $$G$$?
\end{quote}

\section{The key algebraic observation}

Start with the matrix product:

$$
GG^\top
=
\begin{bmatrix}
G_0 & G_1 & \cdots & G_{p-1}
\end{bmatrix}
\begin{bmatrix}
G_0^\top \\
G_1^\top \\
\vdots \\
G_{p-1}^\top
\end{bmatrix}.
$$

Multiplying this block matrix gives

$$
GG^\top
=
G_0G_0^\top
+
G_1G_1^\top
+
\cdots
+
G_{p-1}G_{p-1}^\top.
$$

This is the central reason the left preconditioner is the nice case.

Each rank can compute its own local contribution

$$
P_r = G_rG_r^\top,
$$

and then the final matrix is just

$$
L_{\text{fresh}}
=
\sum_{r=0}^{p-1} P_r.
$$

So we do not need a ring exchange of gradient shards for $$GG^\top$$. Each rank can compute a partial contribution locally, and the ranks only need to sum the partials.

\section{The naive implementation}

The most direct implementation would be:

\begin{enumerate}
    \item Every rank computes its full local contribution:
    $$
    P_r = G_rG_r^\top.
    $$

    \item All ranks sum the full matrix using an all-reduce:
    $$
    P = \sum_r P_r.
    $$

    \item Every rank gets the full $$m \times m$$ matrix.
\end{enumerate}

This is simple, but bad.

Every rank temporarily stores a full matrix:

$$
m \times m.
$$

Every rank also communicates the full matrix. This defeats the point of distributed storage.

So the first improvement is:

\begin{quote}
Do not all-reduce the whole matrix to every rank. Sum the matrix and distribute it directly into the layout needed by the eigensolver.
\end{quote}

\section{The target layout: 2D block-cyclic storage}

Dense eigensolvers such as ScaLAPACK/ELPA-style solvers usually expect the matrix to be distributed across a 2D process grid.

Assume the $$p$$ ranks are arranged as a $$P_r \times P_c$$ grid:

$$
p = P_rP_c.
$$

A rank $$q$$ has process-grid coordinates

$$
\text{grid}_r = q \mathbin{//} P_c,
$$

$$
\text{grid}_c = q \bmod P_c.
$$

Now split the matrix $$L$$ into square blocks of size

$$
b \times b.
$$

Let the global block coordinates be

$$
(i_b, j_b).
$$

The block-cyclic ownership rule is:

$$
\text{owner}(i_b, j_b)
=
(i_b \bmod P_r)P_c
+
(j_b \bmod P_c).
$$

So rank $$q$$ owns block $$(i_b,j_b)$$ exactly when

$$
i_b \bmod P_r = \text{grid}_r,
$$

and

$$
j_b \bmod P_c = \text{grid}_c.
$$

The local block coordinates are

$$
\text{local block row} = i_b \mathbin{//} P_r,
$$

$$
\text{local block col} = j_b \mathbin{//} P_c.
$$

This means each rank stores a dense local matrix of shape

$$
\left(\frac{n_b}{P_r}b\right)
\times
\left(\frac{n_b}{P_c}b\right),
$$

where

$$
n_b = \frac{m}{b}.
$$

\section{From all-reduce to reduce-scatter}

Instead of all-reducing the whole matrix, we want each rank to receive only the pieces of $$L$$ that it owns.

That is exactly the role of a reduce-scatter.

Conceptually, every rank prepares buckets:

\begin{lstlisting}
bucket for rank 0
bucket for rank 1
...
bucket for rank p-1
\end{lstlisting}

Each bucket contains that rank's local partial contribution to the blocks owned by the destination rank.

Then reduce-scatter performs two operations at once:

\begin{enumerate}
    \item It sums corresponding buckets across all ranks.
    \item It gives each rank only its own final bucket.
\end{enumerate}

So rank $$q$$ receives

$$
\sum_{r=0}^{p-1} \text{bucket}_{r \rightarrow q}.
$$

This is the final accumulated piece of $$GG^\top$$ owned by rank $$q$$.

\section{The first streamed version: full horizontal bands}

A better implementation avoids forming the full $$m \times m$$ matrix at once.

Instead, process horizontal panels of block rows.

Let a panel start at block row $$i_0$$ and contain $$c$$ block rows. In row coordinates, the panel is

$$
G[i_0b : (i_0+c)b, :].
$$

A simple streamed version computes

$$
P_{\text{panel}}
=
G_r[i_0b : (i_0+c)b, :]
G_r^\top.
$$

This has shape

$$
(cb) \times m.
$$

Visually, this computes a full horizontal band:

\begin{lstlisting}
        c0 c1 c2 c3 c4 c5 c6 c7
r0      X  X  X  X  X  X  X  X
r1      X  X  X  X  X  X  X  X
r2      X  X  X  X  X  X  X  X
r3      X  X  X  X  X  X  X  X
\end{lstlisting}

This is much better than storing a full $$m \times m$$ temporary. But it still wastes work.

Why?

Because

$$
GG^\top
$$

is symmetric.

So if the downstream eigensolver consumes only the lower triangle, then the upper triangle is redundant.

\section{Using symmetry: why lower triangle matters}

Since

$$
L = GG^\top,
$$

we have

$$
L^\top = L.
$$

So

$$
L_{ij} = L_{ji}.
$$

Therefore, only one triangle is needed. We choose the lower triangle:

$$
i_b \geq j_b.
$$

Visually, the useful blocks are:

\begin{lstlisting}
        c0 c1 c2 c3 c4 c5 c6 c7
r0      X
r1      X  X
r2      X  X  X
r3      X  X  X  X
r4      X  X  X  X  X
r5      X  X  X  X  X  X
r6      X  X  X  X  X  X  X
r7      X  X  X  X  X  X  X  X
\end{lstlisting}

So the full-band streamed version is still doing almost twice the necessary work and communication.

The goal is now:

\begin{quote}
Compute only the lower-relevant part of each panel, pack it into 2D block-cyclic reduce-scatter order, and update only the local lower triangle.
\end{quote}

\section{The fully ragged lower-triangle idea}

The most mathematically minimal approach would pack only the exact lower-triangular tiles.

For example, for a panel containing block rows $$4,5,6,7$$, the exact useful region is:

\begin{lstlisting}
        c0 c1 c2 c3 c4 c5 c6 c7
r4      X  X  X  X  X
r5      X  X  X  X  X  X
r6      X  X  X  X  X  X  X
r7      X  X  X  X  X  X  X  X
\end{lstlisting}

This avoids all redundant upper-triangle blocks.

But it creates an irregular packing problem. Different destination ranks may receive different numbers of tiles. The code becomes more complex, and packing may require many small strided copies.

This is theoretically clean, but it may not be the fastest practical GPU implementation.

\section{The chosen compromise: lower-prefix streaming}

The implementation uses a lower-prefix strategy.

For a panel ending at block row $$i_0+c$$, instead of computing against all rows of $$G$$, we compute only against the prefix:

$$
G_r[0 : (i_0+c)b, :].
$$

So the panel product is

$$
P_{\text{panel}}
=
G_r[i_0b : (i_0+c)b, :]
\;
G_r[0 : (i_0+c)b, :]^\top.
$$

Its shape is

$$
(cb) \times ((i_0+c)b).
$$

This avoids all future-column work.

Visually, for a panel covering block rows $$4,5,6,7$$, the full-band version computes:

\begin{lstlisting}
        c0 c1 c2 c3 c4 c5 c6 c7
r4      X  X  X  X  X  X  X  X
r5      X  X  X  X  X  X  X  X
r6      X  X  X  X  X  X  X  X
r7      X  X  X  X  X  X  X  X
\end{lstlisting}

The lower-prefix version computes:

\begin{lstlisting}
        c0 c1 c2 c3 c4 c5 c6 c7
r4      X  X  X  X  X  X  X  X
r5      X  X  X  X  X  X  X  X
r6      X  X  X  X  X  X  X  X
r7      X  X  X  X  X  X  X  X
\end{lstlisting}

Here the prefix ends at $$c7$$ because the panel ends at $$r7$$. For an earlier panel, say rows $$2,3$$, it computes:

\begin{lstlisting}
        c0 c1 c2 c3 c4 c5 c6 c7
r2      X  X  X  X  .  .  .  .
r3      X  X  X  X  .  .  .  .
\end{lstlisting}

The true lower triangle for rows $$2,3$$ is:

\begin{lstlisting}
        c0 c1 c2 c3 c4 c5 c6 c7
r2      X  X  X  .  .  .  .  .
r3      X  X  X  X  .  .  .  .
\end{lstlisting}

So lower-prefix still has some intra-panel upper waste, but it removes all future-column waste.

This is the key tradeoff.

\begin{center}
\begin{tabular}{>{\raggedright\arraybackslash}p{0.24\textwidth}
                >{\raggedright\arraybackslash}p{0.31\textwidth}
                >{\raggedright\arraybackslash}p{0.31\textwidth}}
\toprule
Method & Benefit & Problem \\
\midrule
Full-band streaming
& Simple and regular
& Computes and communicates too much \\
\midrule
Exact ragged lower-only
& Minimum work and communication
& Irregular packing and harder implementation \\
\midrule
Lower-prefix streaming
& Most lower-triangle savings with regular packing
& Small intra-panel upper waste \\
\bottomrule
\end{tabular}
\end{center}

\section{Why the chunk must align with the process grid}

Let the panel contain $$c$$ block rows.

The implementation requires

$$
c \equiv 0 \pmod{\operatorname{lcm}(P_r,P_c)}.
$$

This is not arbitrary. It ensures that the prefix rectangle can be evenly split across both the process rows and process columns.

The prefix block count is

$$
i_0 + c.
$$

If $$i_0$$ and $$c$$ are multiples of $$\operatorname{lcm}(P_r,P_c)$$, then the prefix is divisible by $$P_c$$, and the panel rows are divisible by $$P_r$$.

This allows the panel to be viewed regularly as

$$
\left[
\frac{c}{P_r},
P_r,
b,
\frac{i_0+c}{P_c},
P_c,
b
\right],
$$

then permuted into destination-rank order:

$$
\left[
P_r,
P_c,
\frac{c}{P_r},
b,
\frac{i_0+c}{P_c},
b
\right].
$$

The first two dimensions, $$P_r$$ and $$P_c$$, determine the destination rank:

$$
\text{dst rank}
=
\text{row owner} \cdot P_c
+
\text{column owner}.
$$

This is why the packing can be a clean rectangular copy rather than a ragged tile loop.

\section{The reduce-scatter step}

After packing, the send buffer is organized as:

\begin{lstlisting}
chunk for rank 0
chunk for rank 1
...
chunk for rank p-1
\end{lstlisting}

Each rank has such a send buffer, containing its local partial contribution.

Then reduce-scatter computes:

$$
\text{recv}_q
=
\sum_{r=0}^{p-1}
\text{send}_{r,q}.
$$

So rank $$q$$ receives its own block-cyclic piece of the current prefix panel, already summed across all column shards.

This is exactly the result we need for the EMA update.

\section{Updating only the local lower triangle}

The function stores $$L_{\text{prev}}$$ as a dense local 2D block-cyclic tensor of shape

$$
\left(\frac{n_b}{P_r}b\right)
\times
\left(\frac{n_b}{P_c}b\right).
$$

However, only entries corresponding to global lower-triangular blocks are valid.

For each local block row in the received panel, the global block row is

$$
i_b = i_0 + \text{grid}_r + kP_r.
$$

This rank owns global column blocks

$$
j_b = \text{grid}_c,
\text{grid}_c + P_c,
\text{grid}_c + 2P_c,
\ldots
$$

Only the blocks satisfying

$$
j_b \leq i_b
$$

belong to the lower triangle.

The number of valid local column blocks is therefore

$$
n_{\text{cols}}
=
\left\lfloor
\frac{i_b - \text{grid}_c}{P_c}
\right\rfloor
+ 1.
$$

Multiplying by the block size gives the number of valid local columns in the dense local tensor:

$$
n_{\text{cols, entries}}
=
b
\left(
\left\lfloor
\frac{i_b - \text{grid}_c}{P_c}
\right\rfloor
+ 1
\right).
$$

The EMA update is then applied only to that lower-triangular prefix:

$$
L_{\text{prev, local}}
\leftarrow
\beta L_{\text{prev, local}}
+
(1-\beta)L_{\text{fresh, local}}.
$$

Upper global blocks are intentionally not updated.

\section{The overlap pipeline}

The function also overlaps computation and communication.

For each panel $$t$$, the work is:

\begin{enumerate}
    \item Compute the lower-prefix GEMM.
    \item Pack the result into reduce-scatter order.
    \item Launch the reduce-scatter asynchronously.
    \item Drain the previous panel's reduce-scatter and apply the EMA update.
\end{enumerate}

The timeline is:

\begin{lstlisting}
time ------------------------------------------------------>

GEMM/pack:     panel 0      panel 1      panel 2      panel 3
comm:              RS 0         RS 1         RS 2         RS 3
EMA update:          upd 0        upd 1        upd 2        upd 3
\end{lstlisting}

The important design choice is:

\begin{quote}
Launch the current panel's reduce-scatter before draining the previous panel.
\end{quote}

This keeps the communication queue fed while the previous result is being unpacked and accumulated.

The implementation uses two send buffers and two receive buffers, alternating by panel index. This allows one collective to use one buffer while the next panel uses the other buffer.

\section{Why the workspace cache matters}

A naive implementation would allocate temporary buffers every time the optimizer updates $$L$$.

That is bad because this function may run repeatedly for many layers and many optimizer steps.

So the production-oriented version caches:

\begin{itemize}[leftmargin=1.5em]
    \item the panel output buffer,
    \item the send buffers,
    \item the receive buffers,
    \item and optionally a contiguous fp32 copy of $$G_{\text{local}}$$.
\end{itemize}

In steady state, the function should not allocate large temporary tensors.

This matters because allocator overhead and memory fragmentation can distort performance measurements and hurt training throughput.

\section{The algorithm in one page}

\begin{enumerate}
    \item Interpret $$G$$ as column-sharded:
    $$
    G =
    \begin{bmatrix}
    G_0 & G_1 & \cdots & G_{p-1}
    \end{bmatrix}.
    $$

    \item Use the identity:
    $$
    GG^\top = \sum_r G_rG_r^\top.
    $$

    \item Split $$L$$ into $$b \times b$$ blocks.

    \item Distribute blocks using a $$P_r \times P_c$$ 2D block-cyclic layout:
    $$
    \text{owner}(i_b,j_b)
    =
    (i_b \bmod P_r)P_c
    +
    (j_b \bmod P_c).
    $$

    \item Stream over panels of block rows.

    \item For each panel, compute only the lower-prefix product:
    $$
    P_{\text{panel}}
    =
    G_r[i_0b:(i_0+c)b,:]
    G_r[0:(i_0+c)b,:]^\top.
    $$

    \item Pack the prefix rectangle into destination-rank order.

    \item Reduce-scatter the packed buffers:
    $$
    \text{recv}_q
    =
    \sum_r
    \text{send}_{r,q}.
    $$

    \item Update only local lower-triangular blocks:
    $$
    L \leftarrow \beta L + (1-\beta)\text{recv}.
    $$

    \item Reuse cached workspace buffers across calls.
\end{enumerate}

\section{What this function is not doing}

It is important to be precise about what this implementation does not do.

\begin{itemize}[leftmargin=1.5em]
    \item It does not compute a full $$m \times m$$ matrix.
    \item It does not all-reduce the full matrix to every rank.
    \item It does not compute the entire upper triangle.
    \item It does not store a fully valid dense local matrix; upper global blocks are intentionally stale.
    \item It does not use a ragged exact lower-only tile packer.
\end{itemize}

Instead, it uses a practical compromise:

\begin{quote}
Lower-prefix GEMMs give most of the lower-triangle savings, while regular rectangular packing keeps the implementation GPU-friendly.
\end{quote}

\section{The core tradeoff}

The full-band implementation is simpler, but wastes almost a factor of two.

The exact ragged lower-only implementation is more theoretically efficient, but harder to pack efficiently on GPUs.

The lower-prefix implementation is the middle ground:

$$
\text{most of the lower-triangle savings}
+
\text{regular block-cyclic packing}
+
\text{good overlap structure}.
$$

This is why it is a strong production-oriented candidate.

\section{What must still be tested}

Before trusting this inside SOAP, the minimum tests are:

\begin{enumerate}
    \item Reconstruct the distributed lower triangle and compare against dense $$GG^\top$$.
    \item Test multiple grids:
    $$
    1 \times 8,\quad 2 \times 4,\quad 4 \times 2,\quad 8 \times 1.
    $$
    \item Test multiple chunk sizes:
    $$
    \operatorname{lcm}(P_r,P_c),\quad
    2\operatorname{lcm}(P_r,P_c),\quad
    4\operatorname{lcm}(P_r,P_c),\quad
    8\operatorname{lcm}(P_r,P_c).
    $$
    \item Measure GEMM time, packing time, reduce-scatter wait time, and EMA update time separately.
    \item Fill upper global blocks with NaN in debug mode and verify no downstream code reads them.
\end{enumerate}

\section{Summary}

The function starts from a simple identity:

$$
GG^\top = \sum_r G_rG_r^\top.
$$

That identity lets each rank compute a local partial contribution from its column shard of $$G$$. Instead of all-reducing a full matrix, the function packs each streamed panel into 2D block-cyclic destination order and uses reduce-scatter to both sum the partials and deliver the right shard to each rank. Then, because $$GG^\top$$ is symmetric, the function avoids computing future-column work by using lower-prefix panels of the form

$$
G[i_0:i_1]G[:i_1]^\top.
$$

Finally, it overlaps communication with later panels and reuses cached workspace buffers so the steady-state path avoids large allocations. The result is a practical, production-oriented implementation of exact distributed $$GG^\top$$ for column-sharded gradients.

\end{document}